\documentclass[10pt, a4paper]{article}
\usepackage[margin=2cm]{geometry}
\usepackage[utf8]{inputenc}
\usepackage[portuguese]{babel}
\usepackage{amsmath}
\usepackage{amssymb}
\usepackage{mathtools}
\DeclareMathOperator{\Var}{Var}
\newcommand{\E}{\mathbb{E}}
\newcommand{\prob}{\mathbb{P}}

\title{\textbf{Lista 2 ME323}}
\author{Julio da Silva Telles RA:281275}

\begin{document}

\maketitle{}

\begin{enumerate}
    \item \begin{enumerate}
        \item Classica
        \item Frequência relativa
        \item Modelagem
    \end{enumerate}
    
    \item Denotando o espaço amostral como $\Omega \text{ e o conjunto de eventos favoraveis de } {\E}$ temos:
    $$\prob(\E) = \frac{n(\E)}{n(\Omega)}$$ tal que $\prob(\E)$ = 37,5%.
    
    \item Dado que $\prob(A) = 0.18$ temos que: $$\prob(A^{c}) = 1 - \prob(A)$$ Portando $\prob(A) = 0.82$.
    
    \item \begin{enumerate} 
        \item Dado que $\prob(N) = 0.7$, $\prob(T) = 0.55$ e $\prob(N \cap T) = 0.4$ temos:
        $$\prob(N \cup T) = \prob(N) + \prob(T) - \prob(N \cap T)$$
        Tal que $\prob(N \cup T) =  0.85$.
        \item Sendo $\prob(N \cup T) = 0.85$ temos:
        $$\prob(-) = 1 - \prob(N \cup T)$$
        Tal que $\prob(-) = 0.15$
    \end{enumerate}
    
    \item Denotando o espaço amostral como $\Omega$ tal que $n(\Omega) = 36$ (PFC) e $\mathbb{S}$ como o conjunto
    de combinações de faces cuja soma é 7, tal que $n(\mathbb{S}) = 6$ temos:
    $$\prob(S) = \frac{n(\mathbb{S})}{n(\Omega)}$$
    Assim temos que $\prob(\mathbb{S}) = 1/6$.
    
    \item Denotando o espaço amostral como $\Omega$ tal que $n(\Omega) = 200$ e $\mathbb{T}$ como o conjunto
    dos casos em que ocorreu travamento $n(\mathbb{S}) = 18$ temos:
    $$\prob(S) = \frac{n(\mathbb{T})}{n(\Omega)}$$
    Assim temos que $\prob(\mathbb{T}) = 0.9$.
    
    \item Dado que $\prob(chuva) = 0.25$ e $\prob(atraso \cap chuva) = 0.15$ temos:
    $$\prob(atraso|chuva) = \frac{\prob(atraso \cap chuva)}{\prob(chuva)}$$
    Tal que $\prob(atraso|chuva) = 0.6$.
    
    \item Tendo como verdade que:
    $$\prob(A \cap B) = \prob(A) \times \prob(B) \Longleftrightarrow \text{A e B são independentes}$$
    Assim verificamos que como $0.06 = 0.3 \times 0.2$ A e B são independetes.
    
    \item Sejam os eventos:
    \begin{itemize}
        \item $A, B, C$: a placa é do fornecedor A, B ou C, respectivamente.
        \item $F$: a placa apresenta falha.
    \end{itemize}

    Pelos dados do problema, temos as probabilidades a priori e as probabilidades condicionais (taxas de falha):
    \begin{itemize}
        \item $\prob(A) = 0,50$; $P(F|A) = 0,02$
        \item $\prob(B) = 0,30$; $P(F|B) = 0,04$
        \item $\prob(C) = 0,20$; $P(F|C) = 0,01$
    \end{itemize}

    Utilizando o \textbf{Teorema da Probabilidade Total}, a probabilidade de uma placa escolhida ao acaso apresentar falha é:
    $$\prob(F) = \prob(A) \cdot \prob(F|A) + \prob(B) \cdot \prob(B) + \prob(C) \cdot \prob(F|C)$$
    Tal que $\prob(F) = 0,024$.

    \item Sejam a probabilidade de cada evento ocorrer:
    \begin{itemize}
        \item $\prob(P) = 0.05$
        \item $\prob(+|P) = 0.92$
        \item $\prob(+|P^{c}) = 0.08$
    \end{itemize}
    Temos então que:
    $$\prob(+) =  \prob(P) \cdot \prob(+|P) + \prob(P^{c}) \cdot \prob(+|P^{c})$$
    $$\prob(+) =  \prob(P) \cdot \prob(+|P) + (1 - \prob(P)) \cdot \prob(+|P^{c})$$
    Assim $\prob(+) = 0.122$ (pelo teorema de Bayes) $\Longrightarrow$ 
    $$\prob(P|+) = \frac{\prob(P) \cdot \prob(+|P)}{\prob(+)}$$
    Tal que $\prob(P|+) = 0.377$.

    \item Dado as probabilidades do exercicio anterior temos que:
    $$\prob(+) =  \prob(P) \cdot \prob(+|P) + \prob(P^{c}) \cdot \prob(+|P^{c})$$
    $$\prob(+) =  \prob(P) \cdot \prob(+|P) + (1 - \prob(P)) \cdot \prob(+|P^{c})$$
    Assim $\prob(+) = 0.122$.

    \item Sim, a afirmação está correta. Dado que $\mathbb{B} = \mathbb{A} \cup \mathbb{B - A}$ temos que:
    $$\prob(B) = \prob(A) + \prob(B - A)$$
    Pelo Axioma da não negatividade temos que $\prob(B - A) \ge 0 \Longrightarrow \prob(A) \ge \prob(B)$.
\end{enumerate}

\end{document}
